\section{Limits, maximal honest formulations, and research directions}\label{sec:limits}

\subsection{The framework is not a completeness theorem for mathematics}

\TAPIR{} is complete for neither all constructive mathematics nor all theorems true in the standard
natural numbers.  It is a proof-producing architecture for finite, primitive-recursive, or explicitly
bounded interfaces.  A theorem may fail to fit because:

\begin{itemize}
\item no primitive-recursive presentation of the needed data is known;
\item the proof requires an unbounded search or a nonconstructive selection that cannot be sealed;
\item an effective modulus, height bound, or separation constant is unavailable;
\item the desired statement is stronger than the finite observations supported by the certificates;
\item the theorem is simply not provable in iPRA.
\end{itemize}

The correct response is to weaken the public statement to a finite-use or residualized theorem, not
to hide the obstacle in an informal compiler assumption.

\subsection{Maximal honest formulations}

Several recurring situations have canonical strongest statements.

\paragraph{Convergence without a modulus.}
State: for the fixed finite list of requested precisions, a supplied certificate answers them, or a
metastable interval contains a stable pair.  Do not state a uniform rate.

\paragraph{Finite solutions without an effective height.}
State: any supplied height bound yields a complete bounded list, and prove the list theorem.  Keep the
height theorem as a residual.

\paragraph{Noncanonical exact data.}
State an invariant observation, equality class, or small cloud.  Add exact naming only under a gap.

\paragraph{Compactness-like arguments.}
State the result under a supplied finite net/cover and modulus.  If the domain-specific theory can
produce them, a later cut closes the theorem.

\paragraph{Classical existence with finite verification.}
If the witness is local, existentially eliminate it from an object-level existence theorem.  If only
a host proof of existence is known, no object theorem has yet been produced.

\subsection{Potential extension: proof-plan reification}

The present plan is a Lean datatype whose interpreter produces \(\Prf\).  A later phase may reify
plans as natural codes and prove a bounded checker correct inside iPRA.  Such reification would be
useful for proof compression, self-verification experiments, or Busy Beaver consistency machines.
It should be undertaken only after the direct typed compiler is checked.

The reified language must be deliberately finite: fixed rule set, bounded syntax, explicit fuel, and
no universal primitive-recursive evaluator.  Its soundness theorem would concern accepted plans below
a bound, not arbitrary encoded Lean programs.

\subsection{Potential extension: optimizer correctness}

Because plans expose structure, one can define host transformations:

\begin{itemize}
\item dead-node elimination;
\item common-subproof sharing;
\item fusion of consecutive substitutions;
\item context minimization;
\item choice between ordinary \(\Prf\) and native-clause backends;
\item trace compression.
\end{itemize}

There are two sound implementation strategies.  The simplest optimizer returns a new indexed plan;
type checking plus \code{Plan.run} suffices.  A more ambitious verified optimizer proves equality or
mutual derivability of plans.  The former should be preferred until performance requires the latter.

\subsection{Potential extension: resource certificates}

Proof generation itself has finite resource questions.  A plan may carry:

\begin{itemize}
\item an upper bound on object proof nodes;
\item maximal context length;
\item maximal substitution arity;
\item number of induction uses;
\item certificate search-space size.
\end{itemize}

A host analyzer can report these metrics, and a reified bounded checker could verify some of them.
This would make proof-term explosion measurable across framework alternatives.

\subsection{Potential extension: front ends}

Domain-specific front ends can compile to the same IR:

\begin{enumerate}
\item a recurrence DSL for sequences and bounded sums;
\item an exact polynomial/matrix DSL using the existing ring tactic;
\item a quantitative-analysis DSL with precision propagation and error budgets;
\item a finite-strategy DSL for bounded alternating formulas;
\item a finite-cover DSL for geometry and topology;
\item a height-list DSL for Diophantine packets.
\end{enumerate}

The front ends may be elaborate and untrusted.  Their outputs remain safe because the common typed
back end must elaborate to \(\Prf\,\PRAAx\).

\subsection{Research question: optimal boundary discovery}

The reference compiler computes residuals syntactically.  A stronger analyzer could use proof search
or interpolation to reduce the boundary set, but logical minimality is generally difficult and may be
undecidable.  A pragmatic direction is to annotate each local node with a small declared boundary set
and let backward slicing remove unused declarations.  This gives predictable, auditable results.

\subsection{Research question: constructive observational quotients}

Stable observer erasure suggests a reusable theory of finite observational descent.  Given a graph of
presentations and compatibility edges, one can bound the discrepancy between any two nodes by a path
sum.  A spanning tree certificate suffices.  Formalizing this generically over coded seminorms or
entourages could unify branch clouds, chart transitions, numerical enclosures, and approximate group
representations.

\subsection{Research question: bridge-aware compilation}

The current unification plan supports open-sequent transport among the main kernels.  A future
compiler could select the cheapest backend per node, then transport every result to ordinary
\(\Prf\,\PRAAx\) at the schedule boundary.  Correctness would rely only on checked bridge theorems.
A cost model could compare native clauses, ordinary natural deduction, and specialized arithmetic
lemmas.

\subsection{Honesty about this artifact}

No Lean executable was installed in the runtime used to build this archive.  Therefore:

\begin{itemize}
\item the two no-placeholder Lean files are carefully matched to inspected APIs but unelaborated;
\item files with \code{sorry} are formalization outlines, not proofs;
\item Python test success says nothing about derivability in iPRA;
\item the only kernel-checked facts relied upon are those already present in the pinned repository,
      as documented in the source audit.
\end{itemize}

This status is part of the mathematical result: it prevents a source sketch from being mistaken for a
verified theorem.
